\section{Planning Loop Trace Metrics for Loop Engineering Evaluation}\label{app:planning-metrics}

We formalize the six planning metrics used in RQ2 for loop engineering evaluation. They compare the plan DAG recorded in the evaluated loop trace with the source recovered human DAG and summarize how planning choices change routing burden and obligation retention pressure (Table~\ref{tab:rq2-grid}). Let $G_A = (V_A, E_A)$ denote the recorded plan DAG and $G_H = (V_H, E_H)$ the human development DAG recovered from source records. Let $V_\cap = V_A \cap V_H$ denote the set of units present in both DAGs.

\subsection{Edge F1}\label{app:edge-f1}
Let $E_\cap = E_A \cap E_H$, and define
\begin{equation}
P_{\text{edge}} = \frac{|E_\cap|}{|E_A|}, \qquad R_{\text{edge}} = \frac{|E_\cap|}{|E_H|}.
\end{equation}
\begin{equation}
\text{Edge\,F1} \;=\; \frac{2\,P_{\text{edge}}\,R_{\text{edge}}}{P_{\text{edge}} + R_{\text{edge}}}.
\end{equation}
Edge F1 measures exact recovery of direct prerequisite edges. Low precision corresponds to unsupported dependencies in the recorded plan, while low recall corresponds to omitted prerequisites that can cause work to be routed before prerequisite state is established.

\subsection{TC-F1}\label{app:tc-f1}
Let $\operatorname{TC}(\cdot)$ denote the transitive closure of a DAG's edge set, and write $T_\cap = \operatorname{TC}(E_A) \cap \operatorname{TC}(E_H)$. Define
\begin{equation}
P_{\text{tc}} = \frac{|T_\cap|}{|\operatorname{TC}(E_A)|}, \qquad R_{\text{tc}} = \frac{|T_\cap|}{|\operatorname{TC}(E_H)|}.
\end{equation}
\begin{equation}
\text{TC\,F1} \;=\; \frac{2\,P_{\text{tc}}\,R_{\text{tc}}}{P_{\text{tc}} + R_{\text{tc}}}.
\end{equation}
TC-F1 is a more permissive measure of prerequisite ordering agreement. It credits a recorded plan that preserves \emph{$u$ before $v$} even when it uses a different intermediate path from the human development DAG, capturing whether broad routing order survives edge level mismatch.

\subsection{Layer $\rho$}\label{app:layer-rho}

Let $\ell_A(v)$ and $\ell_H(v)$ denote the topological layer index of unit $v$ in $G_A$ and $G_H$. Layer~0 contains prerequisite free nodes, and later layers depend on earlier layers.
\begin{equation}
\begin{aligned}
\operatorname{Layer}\rho \;=\; \operatorname{Spearman}\!\bigl(& \{\ell_A(v)\}_{v \in V_\cap}, \\
& \{\ell_H(v)\}_{v \in V_\cap}\bigr).
\end{aligned}
\end{equation}
Layer~$\rho$ records whether the plan places units at similar development \emph{stages} as the human reference, independent of exact edge structure. Low values correspond to the evaluated loop assigning units to different prerequisite layers, which changes when obligations become ready during execution.

\subsection{Critical Path Ratio}\label{app:cp-ratio}
Let $\operatorname{CP}(G)$ denote the length (in nodes) of the longest directed path in $G$.
\begin{equation}
\operatorname{CPR}(G) \;=\; \frac{\operatorname{CP}(G)}{|V|}.
\end{equation}
Values near~$1$ correspond to a nearly linear recorded plan that overserializes the task. Lower values correspond to preservation of parallel branches that can be routed independently. We also report a human normalized variant $\operatorname{CPR}(G_A)/\operatorname{CPR}(G_H)$ in extended results to separate task shape from loop planning behavior.

\subsection{Width Ratio}\label{app:width-ratio}
Let $w(G) = \max_\ell |\{v \in V : \ell(v) = \ell\}|$ denote the maximum topological layer width of $G$.
\begin{equation}
\operatorname{WR} \;=\; \frac{w(G_A)}{w(G_H)}.
\end{equation}
Values close to~$1$ correspond to a recorded plan that matches the human development DAG's concurrent ready frontier width. Values below~$1$ correspond to overserialization, while values above~$1$ are consistent with overparallelization caused by missing prerequisite edges and therefore higher regression obligation pressure.

\subsection{Branch Switching Rate}\label{app:branch-switch}
Assign each unit $v$ a branch identifier $b(v)$ corresponding to its earliest reachable source node in $G_H$ (units with multiple source ancestors are labeled \emph{mixed}). Given the planned or executed implementation order in the recorded loop trace $\sigma = (v_{\pi(1)}, \ldots, v_{\pi(m)})$, let $K = |\{\,i : b(v_{\pi(i)}) \neq b(v_{\pi(i+1)})\,\}|$.
\begin{equation}
\operatorname{BSR} \;=\; \frac{K}{m-1}.
\end{equation}
High values correspond to breadth first frontier coverage across branches, while low values correspond to depth first completion of one branch before switching. BSR is graph conditional and is read relative to the shape of $G_H$, since the same switching pattern can reflect frontier coverage in a wide DAG or unnecessary alternation in a narrow chain.
